\section{Methodology}
\label{sec:methodology}

This section describes the calibration method that selects and weights records, and then the design
of the experiment that evaluates the method as a sampler.

\subsection{Selection and weighting as one optimization}

The calibration matrix $\mathbf{M} \in \mathbb{R}^{m \times n}$ maps $n$ candidate records to $m$
targets. Entry $M_{ji}$ is the contribution of candidate $i$ to target $j$: an indicator or count
for count targets, and a simulated tax or benefit value for dollar targets. The matrix is built from
\policyengine{} simulations evaluated under the policy rules that apply to each record's assigned
geography, and is stored in sparse form.

Given $\mathbf{M}$, the target vector $\mathbf{t}$, and a weight vector $\mathbf{w} \in \mathbb{R}^n$,
the optimizer minimizes
\begin{equation}
    \mathcal{L}(\mathbf{w}, \boldsymbol{\alpha}) =
    \frac{1}{m} \sum_{j=1}^{m}
    \min\!\left(\left|\frac{\hat{t}_j - t_j}{s_j}\right|, c\right) +
    \lambda_{L_0} \sum_{i=1}^{n} \bar{z}_i +
    \lambda_{L_2} \frac{1}{n} \sum_{i=1}^{n} \left(\frac{w_i}{w_{0,i}}\right)^2,
    \label{eq:loss}
\end{equation}
where
\begin{equation}
    \hat{t}_j = \sum_{i=1}^{n} M_{ji} w_i z_i.
    \label{eq:estimate}
\end{equation}
The first term is a calibration loss: the mean absolute relative error of the estimates against the
targets, where $s_j$ is the scale of target $j$ (its own magnitude $|t_j|$ for a pure relative error)
and $c$ caps the contribution of any single hard-to-fit target, so one badly conditioned target cannot
dominate the gradient. The implementation also admits optional per-target loss weights, which we leave
uniform here. The second term is the $L_0$ penalty: $\bar{z}_i$ is the expected activation of gate $i$
under the Hard Concrete relaxation \citep{louizos2018}, so $\sum_i \bar{z}_i$ is the expected number of
open gates --- the expected retained-record count --- and $\lambda_{L_0}$ sets how aggressively the
optimizer prunes. The third term is a soft concentration penalty: $\lambda_{L_2}$ multiplies the mean
squared ratio of each fitted weight to its initial weight $w_{0,i}$, which discourages the optimizer from
carrying the fit on a few heavily inflated records. It applies to the informed $L_0$ condition and is
left at zero for the baselines. The weights are parameterized in log space so the fitted weights stay
positive. This is the production optimizer that jointly fits positive log-weights and gates, rather than
a configuration that applies a gate on top of a separately fitted weight.

Read as a sampler, the three loss terms have distinct jobs, and a hard per-record constraint backstops
them. The calibration loss keeps the retained sample faithful to the targets, and because it is
relative, a one per cent error on a small target and on a large target count the same. The $L_0$
penalty sets the size of the retained sample. An $L_0$ penalty on its own, however, can reach a sparse
solution that is still unusable, because the optimizer can match the targets by placing very large
weights on a few retained records. The method governs that concentration through the ratio of each
fitted weight to its initial weight, $w_i / w_{0,i}$, with two complementary controls. The $L_2$ term
is a soft penalty on the mean square of that ratio: it discourages large inflations smoothly, and it is
the knob we sweep to trace the accuracy--concentration trade-off. \texttt{max\_weight\_ratio} is a hard
per-record cap on the same ratio: no fitted weight may exceed a fixed multiple of that record's initial
weight, enforced by clamping the log-weights after every optimizer step. It bounds the worst case rather
than the average. The two act at different points: the $L_2$ penalty shapes the whole weight distribution
during optimization, while the cap constrains the largest weight. Size and concentration are therefore
separate knobs: $\lambda_{L_0}$ sets how many records remain, and the $L_2$ penalty and
\texttt{max\_weight\_ratio} set how concentrated their weights may become. Together they let the pipeline
trade off calibration fit, sample size, and weight stability. The headline runs leave the
cap off and fix $\lambda_{L_2}$ at \tbc[the value chosen from the operability sweep of
Section~\ref{sec:results}]; tightening either bounds how much population any single record can carry, at
some cost to calibration fit. Appendix~\ref{app:hyperparameters} lists the full configuration, and
Algorithm~\ref{alg:l0} gives the optimization loop.

The gates and the weights are learned together: the gate decides whether a record is in the sample,
and the weight decides how much it counts. Because both are fit against the same calibration loss,
selection is informed by the targets: a record survives when keeping it helps match a target that
would otherwise be missed. This is the sense in which the method is a target-informed sampler rather
than a reweighting of a fixed random draw.

\subsection{From fitted weights to a dataset}

After optimization the gates are evaluated deterministically, so the retained-record count is read
from the solution rather than chosen by a post-hoc cutoff. The retained records and their weights are
assembled into a dataset for \policyengine{}. Appendix~\ref{app:assembly} summarizes the assembly
step, which matters operationally but is not the object of the experiment.

\subsection{Sampling experiment design}
\label{sec:experiment_design}

The experiment evaluates whether informed selection produces a better dataset than the alternative
samplers at the same record budget. From one candidate universe and one target set, we compare informed
$L_0$ selection against two matched-budget baselines, which span where calibration enters the reduction
(Section~\ref{sec:reduction_methods}):

\begin{itemize}
    \item \textbf{Informed $L_0$ sampling with Hard Concrete gates.} The method above selects records
    and fits weights jointly at a chosen budget, set through $\lambda_{L_0}$.
    \item \textbf{Random sampling with gradient-descent reweighting.} A random subset of the candidate
    universe of the same size is drawn, and its weights are fit to the same targets by gradient
    descent, using the same calibration loss without the gate term. This is the reduce-first,
    calibrate-after baseline.
    \item \textbf{Survey-weight sampling (Hansen--Hurwitz).} The full universe is first reweighted to the
    targets by gradient descent, and records are then drawn with probability proportional to those weights,
    \emph{with replacement}; each of the $n$ draws receives an equal share $W/n$ of the total weight $W$,
    and repeated draws accumulate on the same record. This is the unbiased Hansen--Hurwitz integerisation
    \citep{hansen1943} of the dense weights (it preserves their expectation at every budget), and is
    the calibrate-then-reduce baseline.
\end{itemize}

Combinatorial optimisation, the discrete gradient-free sibling of informed $L_0$ selection reviewed in
Section~\ref{sec:reduction_methods}, is not implemented in this study; the comparison here is the
three-way benchmark above.

Holding the candidate universe and the target set fixed isolates the effect of how records are
chosen. We sweep the retained-record budget across 2{,}000, 5{,}000, 10{,}000, 20{,}000, and 40{,}000 records,
from aggressive compression of the candidate universe up to a large fraction of it, and run each budget
over five calibration seeds (seeds $0$--$4$). Informed $L_0$ sets the achieved budget
at each grid point through its $\lambda_{L_0}$ bisection, and the baselines are matched to the same
retained count.

\subsection{Out-of-sample target evaluation}
\label{sec:oos}

A sampler can match the targets it is fit to and still generalize poorly. To test generalization, the
design separates the targets used to fit the weights from the targets used to score the result, and
holds out \emph{whole target families} rather than a random subset of targets. The reason is structural:
targets within a family nest (a national total is the sum of its state cells), so a random target split
leaks: a held-out cell is nearly determined by its retained siblings, which would overstate
generalization. Holding out a whole family removes that leakage.

The fixed holdout comprises the \texttt{cms\_medicaid}, \texttt{usda\_snap}, and \texttt{state\_income\_tax}
families, about \tbc[out-of-sample target count, $\approx$200] targets spanning three domains (Medicaid
enrollment and spending, SNAP, and state income tax). These families are held out of every method's fit
and scored only out-of-sample. Each held-out family keeps a \emph{sibling} family in the fit (ACA and
Medicare enrollment alongside Medicaid; the federal IRS SOI aggregates alongside state income tax), so
the test asks each method to reproduce a held-out family while a related family remains in the fit,
rather than orphaning a domain entirely. The \texttt{cbo} family is additionally held out as a \populace{}
validation-only family (forecast and aging projections used as diagnostics rather than contemporaneous
targets), so it is never fit. Of the \tbc[total target count] targets in the system,
\tbc[fit target count] are used to fit the weights.

As a stricter robustness check we also rotate whole families through the holdout in a family-grouped
five-fold design: the families are dealt into five folds balanced by target count, and each fold is held
out in turn. Because the folds hold out structurally different families rather than independent draws, we
report them per fold rather than as a single inferential test (Section~\ref{sec:results-generalization});
the rotation surfaces in particular the fold that holds out the dominant \texttt{irs\_soi} family.

\subsection{Evaluation metrics}
\label{sec:metrics}

Each run reports, for the scored targets in-sample and out-of-sample:
\begin{itemize}
    \item the \emph{median} absolute relative error as the headline accuracy measure, with the mean and
    maximum reported alongside but flagged as tail-sensitive;
    \item both a \emph{micro}-average over all scored targets (which the largest family, IRS SOI,
    dominates) and a family \emph{macro}-average (the mean of the per-family errors, giving each family
    equal weight regardless of how many targets it contributes);
    \item the effective sample size $\mathrm{ESS} = \left(\sum_i w_i\right)^2 / \sum_i w_i^2$ and the
    implied design effect, reported as a primary result rather than a footnote, because a sampler can fit
    the targets while concentrating population mass on a few records;
    \item the retained-record count, the largest fitted weights, and runtime.
\end{itemize}
Errors are reported by geographic level and by target family as well as in aggregate, because a sampler
can match national totals while missing local ones, or match one source while missing another.

A small number of targets have denominators so small that a single record can swing their relative error
past 100 per cent. We identify these structurally rather than by an arbitrary cutoff: target $j$ is
denominator-degenerate when its value is smaller than its \emph{identifiability floor}
$\max_i |M_{ji}|\, w_{0,i}$, the largest contribution any single record makes at its initial weight. For
such a target the absolute relative error reflects integer placement noise rather than calibration
quality. Rather than winsorize or trim the error distribution, we \emph{name} these targets in a one-time
audit and report a targeted-removal sensitivity (the mean recomputed with exactly those named targets
dropped) beside the full mean, so the effect of the degenerate targets is visible and attributable.
Target counts are read from each run's manifest.
